\scanpage{107}
\mspage{103}
% archive sheet number: 106

La question se pose dans quelle mesure ces équivalences \deleted{sont valables} \unsure{gardent} un sens et restent valables pour une correspondance entre disjoncteurs généraux. Faute d'avoir une \unsure{décomposition} du type $\alpha_*=f'_*\,{}^{t}f_*$, \deleted{il y a} l'interprétation ``naïve'' ou ``\unsure{directe}'' en termes des propriétés des $f$, $f'$, \deleted{\unsure{va}} semble destinée à disparaître. En fait, ce sont les équivalents algébriques ``\unsure{rappelés}'' tantôt, qui sont en \unsure{premier lieu|première ligne}.

% margin note written sideways
% slanted
\margin{Voir \illegible{1} \unsure{$2^{\mathrm{me}}$ pt} \unsure{CF} \uline{VIII} \unsure{exemple} \unsure{non général} (\unsure{cont.} \unsure{top.}) \unsure{72} \unsure{$\Downarrow|\downarrow$} $\ldots$}
Malheureusement, je ne sais prouver les équivalences de i) à (iii') de A) \inserted{(\unsure{ou} que i) \unsure{implique} \unsure{(ii)} (iii))}, que moyennant des hypothèses draconiennes, proches déjà de l'hypothèse booléenne (\unsure{cf} \unsure{donc} \unsure{aux.} \unsure{Axiomes|Nécessités} des \unsure{repères} à propos des supports). Peut-être est-ce là encore des faux problèmes, vu que pour mes propos géométriques, on ne travaillera jamais qu'avec des supports très particuliers (essentiellement, les supports constructibles) : et on n'aura à prendre des $\operatorname{Inf}$ et $\unsure{=}$ des $\operatorname{Sup}$ que sous des conditions idoines de ``\unsure{commensurabilité}''. Ce sont les cas où, précisément, les arguments
